% ============================================================
% kbs_introduction.tex — KBS journal Introduction (KE-emphasized draft).
% Bibliography keys: kbs_ke_references.bib (new, verified) +
% paper/references.bib (conference refs carried over).
% Numbers marked [E1/E2/E4 update] are placeholders pending the journal
% experiments in plan/extended_journal_paper.md; the conference-scale
% values are shown for continuity.
% ============================================================

\section{Introduction}
\label{sec:intro}

Knowledge-based systems inherit their competence from the knowledge they
acquire, and the cost of that acquisition is the field's oldest
bottleneck~\cite{feigenbaum_knowledge_engineering}. Whether the knowledge
encodes clinical practice, engineering diagnosis, or the subject of a
document, it must be elicited, represented in a vocabulary the system can
use, and validated before it can be trusted to drive inference. The
classical knowledge-engineering cycle makes this explicit: acquire, model,
validate, and deploy, with each step auditable~\cite{studer_knowledge_engineering}.
What has changed is the acquisition channel. Large language models (LLMs)
can now propose structured knowledge from minimal input at negligible
marginal cost, and a fast-growing literature shows them populating
ontologies and knowledge graphs from text~\cite{ciatto_llm_oracle_kbs,
pan_llm_kg_roadmap, ji_kg_survey, xu_llm_gen_ie}. Yet the knowledge
engineering of LLM-acquired knowledge has not kept pace with the acquisition
itself. The literature measures how often an LLM matches a single reference
answer, or how well it extracts triples, but rarely asks the three questions
that decide whether machine-acquired knowledge should be deployed at all:
\textit{Is it valid against professional consensus rather than one
reference? Which representation of it does the consuming system actually
use? And what is it worth given how complete the knowledge base already
is?}

We answer these questions in a domain where all three can be measured
against professional ground truth: bibliographic subject knowledge. A library
catalog describes what each work is about through two canonical knowledge
representations: subject headings---an open, pre-coordinated thesaurus
vocabulary (Library of Congress Subject Headings, LCSH)---and
classification---a closed, hierarchical taxonomy (Dewey Decimal
Classification, DDC). Both are produced professionally, both are validated
against national authority files, and both are consumed by retrieval systems
whose quality is measurable. The domain is therefore a controlled laboratory
for a general knowledge-engineering question: \emph{when an LLM acquires
subject knowledge, is that knowledge valid, which representation carries its
utility, and where does it pay off?}

Prior work on LLMs and bibliographic knowledge stops one step short of this
question on each side. Cataloging studies show that chatbots and LLMs can
draft headings and call numbers with varying agreement against a catalog
record~\cite{holstrom_llm_cataloging, characterising_ai_cataloguing,
ital_genai_cataloging, dobreski_chatbots_cataloging}, but evaluation is
single-model, single-source, and stops at agreement---the very weakness the
knowledge-validation literature warns
about~\cite{okeefe_kbs_validation}. On the retrieval side, metadata-aware
and retrieval-augmented systems demonstrate that structured metadata
improves discovery~\cite{bm25f, folkrag, bergen, lund_rag_libraries}, but
they consume metadata as given; they never ask whether machine-acquired
metadata is good enough to consume. The two literatures do not meet: one
measures acquisition without a downstream system, the other uses a system
without measuring acquisition. In knowledge-engineering terms, neither
completes the acquire--validate--utilize cycle that
Studer~et~al.~\cite{studer_knowledge_engineering} describe as the discipline
of the field, and neither prices the machine-acquired knowledge against the
expert knowledge it supplements~\cite{gruber_portable_ontologies}.

This paper closes that gap. We treat an LLM as an automated
knowledge-acquisition component operating in the retrospective-conversion
setting---proposing subject headings and classification from a sparse record
(title, author, year)---and carry its output through the full cycle:
validation against multi-level professional gold with live authority
checking, and utility measurement in the consuming system, a
metadata-aware retrieval index. Using a corpus of \TODO{E1: 1,300} books
with professionally validated ground truth and a released two-sided
benchmark, we contribute:

\begin{enumerate}
\item \textbf{A complete acquisition-to-utilization evaluation protocol.}
  LLM-acquired subject knowledge is validated against professional gold at
  three levels (exact, semantic-equivalence, acceptable), checked against
  the Library of Congress authority file, and measured by its effect on the
  consuming retrieval system---so acquisition quality is reported as
  \emph{what it does to the system that uses it}, not as isolated
  agreement.
\item \textbf{Representation-dependent utility.} The same acquisition
  channel is evaluated on two knowledge representations. The open,
  pre-coordinated thesaurus representation carries the entire downstream
  retrieval gain (subject headings: $+0.382$ topical nDCG@10, 95\% CI
  $[+0.333,+0.431]$ at conference scale \TODO{E1 update}); the closed,
  hierarchical classification contributes none ($-0.005$, CI
  $[-0.011,+0.000]$). This is a result about \emph{which representation of
  machine-acquired knowledge a retrieval task consumes}, with implications
  that outlive the library setting.
\item \textbf{Completeness-dependent value.} The utility of
  machine-acquired knowledge is largest exactly where the knowledge base
  has gaps: retrieval gain from acquisition falls monotonically as existing
  subject coverage rises, but never reaches zero. Machine-acquired knowledge
  is a gap-filling channel, not a replacement for expert knowledge---and the
  answer to ``which knowledge to acquire first'' is read off the coverage
  curve.
\item \textbf{Validated acquisition accuracy.} Scored against professional
  gold \emph{and} an inter-cataloger agreement ceiling
  \TODO{E2: panel}, with a full knowledge-quality error taxonomy
  (invented, over-specific, under-specific, authority-violating, and
  valid-but-unmatched headings) and a decomposition of the residual
  \TODO{E2}.
\item \textbf{Cost of machine- vs.\ expert-acquired knowledge.} Expert time
  to acquire the same knowledge from scratch versus verifying the LLM
  output \TODO{E4: panel}, so the utility numbers can be read against what
  acquisition actually costs.
\end{enumerate}

Our results are domain-specific in their data and general in their
mechanism. The measured conclusion is that an LLM can act as a
\textit{gap-filling knowledge-acquisition component}: its subject knowledge
is validated as operating near the human agreement ceiling
\TODO{E2: panel + calibration}, its utility is carried by the
representation the consuming system actually uses, and its value
concentrates in the low-completeness regime where expert knowledge is
missing. For knowledge-based systems beyond libraries, the transferable
claim is the method and the two regularities---representation dependence and
completeness dependence---rather than any library-specific number.

The remainder of the paper is organized as follows. Section~\ref{sec:rel}
positions this work within knowledge engineering and metadata-aware
retrieval. Section~\ref{sec:method} describes the acquisition component, the
two knowledge representations, and the measurement apparatus.
Section~\ref{sec:data} details the corpus and validated ground truth.
Sections~\ref{sec:valid}--\ref{sec:util} report acquisition validity and
downstream utility respectively. Section~\ref{sec:cost} reports acquisition
cost, Section~\ref{sec:synth} synthesizes the deployment analysis, and
Section~\ref{sec:disc} discusses limitations and generality.
